\chapter{Constrained QUBO Formulation}
\label{ch:qubo}

\section{Overview}

The pruning stage chooses, in every cycle, which edges of $G_t$ to keep. The
choice is combinatorial: it must maximise the usefulness of the kept edges
while respecting per-rover degree caps, a global cost budget and the
connectivity of a set of required rovers. This chapter writes that problem as
a single QUBO, explains every term, proves the two properties that make the
encoding correct, works through a five-rover example by hand and reports how
the implementation was verified.

\section{Decision Variables and Notation}

Let $H\subseteq E_t$ be the \emph{hard-critical} edges, which are always kept
(for example a relay rover's only link), and $O=E_t\setminus H$ the optional
edges. Let $R$ be the required rovers with root $r\in R$ and $m=|R|-1$, let
$D_i$ be the degree cap of rover $i$, and let $K$ be the global budget in
integer cost units. Writing $\deg_H(i)$ for the number of hard edges at $i$
and $c_H$ for their total cost, the remaining degree capacity of rover $i$ is
$D_i-\deg_H(i)$ and the remaining budget is $K-c_H$.
Table~\ref{tab:qubo:vars} lists the four families of binary variables.

\begin{table}[htbp]
\centering
\caption{Variable families of the QUBO. Every integer is binary-encoded.}
\label{tab:qubo:vars}
\small
\begin{tabularx}{\linewidth}{@{}l X l@{}}
\toprule
Variable & Meaning & Number of bits \\
\midrule
$x_e,\ e\in O$ & keep (1) or prune (0) optional edge $e$ & $|O|$ \\
$s_i,\ i\in V$ & degree slack, $0\le s_i\le D_i-\deg_H(i)$ & $\sum_i\lceil\log_2(D_i-\deg_H(i)+1)\rceil$ \\
$s_C$ & budget slack, $0\le s_C\le K-c_H$ & $\lceil\log_2(K-c_H+1)\rceil$ \\
$f_{uv},f_{vu},q_e,\ e\in E_t$ & two directed flows and a capacity slack per edge, each in $[0,m]$ & $3|E_t|\lceil\log_2(m+1)\rceil$ \\
\bottomrule
\end{tabularx}
\end{table}

For notational convenience $x_e:=1$ for $e\in H$. The vector $z$ stacks all
bits, and $N$ denotes its length.

\section{The Energy Function}

The complete energy is
\begin{align}
H(z) ={}& \sum_{e\in O}(-U_e+\lambda)\,x_e
 + \eta\sum_{e\in O}\bigl(x_e-x_e^{t-1}\bigr)^2 \notag\\
&+ \rho_d\sum_{i\in V}\Bigl(\sum_{e\in\delta_O(i)}x_e+s_i-\bigl(D_i-\deg_H(i)\bigr)\Bigr)^2 \notag\\
&+ \rho_C\Bigl(\sum_{e\in O}c_ex_e+s_C-(K-c_H)\Bigr)^2 \notag\\
&+ \rho_F\sum_{v\in V}\bigl(\mathrm{in}_v-\mathrm{out}_v-b_v\bigr)^2 \notag\\
&+ \rho_{FC}\sum_{e=(u,v)\in E_t}\bigl(f_{uv}+f_{vu}+q_e-m\,x_e\bigr)^2,
\label{eq:qubo:energy}
\end{align}
where $\delta_O(i)$ is the set of optional edges at rover $i$; the total flows
entering and leaving rover $v$ are $\mathrm{in}_v$ and $\mathrm{out}_v$; and
the supplies are $b_r=-m$, $b_v=1$ for $v\in R\setminus\{r\}$ and
$b_v=0$ otherwise. Each squared bracket is a polynomial of degree two in the
bits, so the whole expression expands into the form $z^{\top}Qz+\mathrm{const}$.
Table~\ref{tab:qubo:consts} collects the constants.

\begin{table}[htbp]
\centering
\caption{Constants of the QUBO used throughout the report.}
\label{tab:qubo:consts}
\small
\begin{tabularx}{\linewidth}{@{}l X r@{}}
\toprule
Constant & Role & Value \\
\midrule
$\lambda$ & sparsity penalty per kept edge & 0.08 \\
$\eta$ & stability penalty per changed decision & 0.15 \\
$\rho_d,\ \rho_C$ & degree and budget penalty weights & 100 \\
$\rho_F,\ \rho_{FC}$ & flow conservation and flow capacity weights & 500 \\
cost scale & quantisation of $\bar c_e$ into budget units & 10 \\
utility weights & trust, link, battery, proximity, freshness, task & 0.25, 0.25, 0.10, 0.10, 0.15, 0.15 \\
\bottomrule
\end{tabularx}
\end{table}

\section{The Soft Terms}

The coefficient $-U_e+\lambda$ rewards an edge for its utility and charges a
constant $\lambda$ for keeping it, so edges with $U_e<\lambda$ are never worth
keeping alone and the constant pressure drives solutions towards sparse
graphs. The stability term charges $\eta$ for every decision that changes
relative to the previous cycle; since $x_e^2=x_e$ it is linear,
$\eta\bigl(x_e(1-2x_e^{t-1})+x_e^{t-1}\bigr)$, and it stops the kept topology
from flapping between nearly equal alternatives, which would confuse both the
GNN and the trust model.

\section{Encoding Inequalities with Slack}

Degree caps and the budget are inequalities. The standard device turns
$\sum_k a_kx_k\le B$ into the equality $\sum_k a_kx_k+s=B$ with an integer
slack $0\le s\le B$, writes $s=\sum_j 2^js_j$ in bits, and penalises the
squared residual. Expanding a single squared constraint
$(a^{\top}y-B)^2$ over a bit vector $y$ and using $y_k^2=y_k$ gives
\begin{equation}
(a^{\top}y-B)^2=\sum_k\bigl(a_k^2-2Ba_k\bigr)y_k+2\sum_{k<l}a_ka_ly_ky_l+B^2,
\label{eq:qubo:expand}
\end{equation}
so a constraint contributes $a_k^2-2Ba_k$ to the diagonal entry of each of its
bits, $a_ka_l$ to each symmetric off-diagonal pair, and $B^2$ to the
constant.

\begin{example}[One degree constraint]
Take a rover with cap $D_i=2$, no hard edges and three optional edges
$x_1,x_2,x_3$. The slack $s_i\in\{0,1,2,3\}$ needs two bits, $s_i=s_0+2s_1$,
so the coefficient vector over $(x_1,x_2,x_3,s_0,s_1)$ is $a=(1,1,1,1,2)$ and
$B=2$. Equation~\eqref{eq:qubo:expand} yields the entries of
Table~\ref{tab:qubo:degree}, all multiplied by $\rho_d$.
\end{example}

\begin{table}[htbp]
\centering
\caption{Coefficients generated by one degree constraint with $D_i=2$
(before multiplication by $\rho_d$).}
\label{tab:qubo:degree}
\small
\begin{tabular}{@{}l r l@{}}
\toprule
Bit & Linear term $a_k^2-2Ba_k$ & Pairwise term $2a_ka_l$ \\
\midrule
$x_1,x_2,x_3$ & $-3$ each & 2 with each other and with $s_0$; 4 with $s_1$ \\
$s_0$ & $-3$ & 2 with each $x$; 4 with $s_1$ \\
$s_1$ & $-4$ & 4 with each $x$ and with $s_0$ \\
constant & $B^2=4$ & \\
\bottomrule
\end{tabular}
\end{table}

Three checks confirm the intended behaviour. Keeping two edges with both
slack bits zero gives $(2-2)^2=0$. Keeping one edge with $s_0=1$ also gives
$0$. Keeping all three edges gives at least $(3-2)^2=1$ for every slack
setting, because slack cannot be negative; the violation therefore costs at
least $\rho_d=100$, while the third edge can gain at most about $1.2$.

\section{The Budget Constraint}

The budget term has the same structure, with the integer costs $c_e$ as
coefficients and $K-c_H$ as right-hand side, and its slack needs only
$\lceil\log_2(K-c_H+1)\rceil$ bits. A budget below the cost of every connected
subgraph is \emph{infeasible} together with connectivity, and no penalty
weighting can change that; Section~\ref{sec:sol:ladder} shows how the
production solver detects and records this case, and
Section~\ref{sec:res:audit} shows that the project's own earlier experiments
fell into exactly this trap.

\section{Exact Connectivity by Single-Commodity Flow}

The two flow terms enforce connectivity of the required set inside the
optimisation itself, rather than through a check applied afterwards.

\begin{proposition}[Flow encodes connectivity]
Let $x$ be a selection of edges. There exist integer flows and slacks that
make both flow terms of Eq.~\eqref{eq:qubo:energy} zero if and only if every
required rover is connected to the root $r$ in the subgraph of selected
edges.
\end{proposition}

\begin{proof}
($\Leftarrow$) If every required rover is connected to $r$, choose a spanning
tree of the root's component and route one unit from $r$ to each required
rover along its unique tree path. Every arc then carries at most $m$ units,
conservation holds at every node with the given supplies, and each selected
edge satisfies $f_{uv}+f_{vu}\le m$, so a slack $q_e\in[0,m]$ closes the
capacity equality. Unselected edges carry no flow and take $q_e=0$.
($\Rightarrow$) If both terms are zero, conservation says that $m$ units leave
$r$ and one unit is absorbed at each other required rover, while capacity
forces $f_{uv}=f_{vu}=0$ on every unselected edge. Decomposing the flow into
paths, each required rover receives its unit along a path of selected edges
from $r$, so it is connected to $r$.
\end{proof}



\section{Why the Penalties Dominate}

\begin{proposition}[No violation is profitable]
With $\rho_d=\rho_C=100$, $\rho_F=\rho_{FC}=500$ and soft coefficients of
magnitude at most about $1.2$ per edge, every exact minimiser of
Eq.~\eqref{eq:qubo:energy} satisfies all constraints whenever a feasible
selection exists.
\end{proposition}

\begin{proof}[Argument]
A violated degree or budget constraint contributes at least $\rho=100$,
because its residual is a non-zero integer and enters squared, while each unit
of violation unlocks at most one edge worth less than $1+\lambda+\eta<1.25$.
Flow violations are penalised quadratically against an at most linear gain.
\end{proof}

The argument is confirmed exhaustively in Section~\ref{sec:qubo:verify}.
It also explains the difficulty of Chapter~\ref{ch:solvers}: the same walls
that make violations unprofitable make them unreachable for a sampler that
must climb over them.

\section{A Worked Example on Five Rovers}
\label{sec:qubo:worked}

The whole formulation is visible on an instance small enough to solve by
hand. There are five rovers: the root $0$ and rovers $1$ to $4$. Seven links
are in range. The link $0$--$1$ is hard-critical, being rover~1's only
tether. Rovers $\{0,2,3,4\}$ are required, so rover~1 is a relay that may be
cut off. The degree caps are $(3,2,3,2,2)$ for rovers $0$ to $4$ and the
budget is $K=12$ units. Table~\ref{tab:qubo:worked} gives the utility, the
integer cost and the previous decision of every link.

\begin{table}[htbp]
\centering
\caption{The five-rover worked example: utility $U_e$ from
Eq.~\eqref{eq:arch:utility}, integer cost $c_e$ and previous-cycle decision
$x_e^{t-1}$.}
\label{tab:qubo:worked}
\small
\begin{tabular}{@{}l r r r@{}}
\toprule
Link & $U_e$ & $c_e$ & $x_e^{t-1}$ \\
\midrule
$0$--$1$ (hard) & 0.810 & 2 & 1 \\
$0$--$2$ & 0.900 & 1 & 1 \\
$0$--$3$ & 0.405 & 8 & 0 \\
$1$--$2$ & 0.705 & 3 & 1 \\
$2$--$3$ & 0.750 & 4 & 1 \\
$2$--$4$ & 0.790 & 2 & 0 \\
$3$--$4$ & 0.555 & 6 & 1 \\
\bottomrule
\end{tabular}
\end{table}

\subsubsection{Size} There are six optional edges. With $m=3$ every flow
integer needs two bits, so the flow family has $3\times7\times2=42$ bits; the
degree slacks need $2+1+2+2+2=9$ bits and the budget slack, with $K-c_H=10$,
needs 4. The QUBO therefore has $6+42+9+4=61$ binary variables, even though
the decision itself has only $2^6=64$ possibilities.

\subsubsection{Enumeration} All 64 optional-edge subsets were enumerated.
Only \textbf{3} are feasible. The optimum keeps $\{0\text{--}1,0\text{--}2,
2\text{--}3,2\text{--}4\}$ at a cost of 9 units, with energy
\[
\underbrace{(-0.82)+(-0.67)+(-0.71)}_{\text{utility and sparsity}}
+\underbrace{0.15\times3}_{\text{stability}}=-1.750,
\]
where the three stability charges come from dropping $1$--$2$ and $3$--$4$
and adding $2$--$4$. The other two feasible subsets both score $-1.555$.

\subsubsection{The trap} Thirty of the 61 infeasible subsets violate only
connectivity. The best of them, $\{0\text{--}1,0\text{--}2,1\text{--}2,
2\text{--}3\}$, has a \emph{lower} soft energy of $-1.965$ because it keeps
the attractive link $1$--$2$ and pays only one stability charge, but it cuts
rover~4 off from the root. A pruner that judged edges by soft score alone, or
checked connectivity only afterwards, would prefer it by $0.215$. In the QUBO
the flow terms add at least 500, so its true energy is above $+498$. The
cheapest degree or budget violator sits at $+97.475$: the barrier between it
and the optimum is $99.2$, almost a hundred times the width of the whole
feasible band ($[-1.750,-1.555]$). Fig.~\ref{fig:qubo:ladder} draws this
energy ladder.

\begin{figure}[htbp]
\centering
\begin{tikzpicture}[
  lvl/.style={thick},
  lab/.style={font=\small, anchor=west},
]
\draw[-{Stealth}] (0,0) -- (0,6.2) node[above, font=\small] {energy (not to scale)};
\draw[lvl, green!50!black] (0.3,1.0) -- (2.3,1.0); \node[lab] at (2.4,1.0) {optimum, feasible: $-1.750$};
\draw[lvl, green!50!black] (0.3,1.6) -- (2.3,1.6); \node[lab] at (2.4,1.6) {two other feasible subsets: $-1.555$};
\draw[lvl, red!70!black, dashed] (0.3,0.35) -- (2.3,0.35); \node[lab, red!70!black] at (2.4,0.35) {connectivity trap, soft energy only: $-1.965$};
\draw[lvl, orange!80!black] (0.3,3.6) -- (2.3,3.6); \node[lab] at (2.4,3.6) {best degree/budget violator: $+97.475$};
\draw[lvl, red!70!black] (0.3,5.6) -- (2.3,5.6); \node[lab] at (2.4,5.6) {same trap with flow penalty: $>+498$};
\draw[-{Stealth}, red!70!black] (1.3,0.45) -- (1.3,5.5);
\draw[decorate, decoration={brace, amplitude=5pt, mirror}] (-0.15,3.6) -- (-0.15,1.0) node[midway, left=6pt, font=\small, align=right] {barrier\\$99.2$};
\end{tikzpicture}
\caption[Energy ladder of the worked example.]{Energy ladder of the worked example. The connectivity trap looks
better than the optimum to a pruner that ignores connectivity; the flow
penalty lifts it far above every feasible subset.}
\label{fig:qubo:ladder}
\end{figure}

\subsubsection{Cross-checks} The Q-matrix energy of the optimum, computed as
$z^{\top}Qz+\mathrm{const}$ on the feasible bit vector, equals $-1.750$ and
matches enumeration. The production annealer of Chapter~\ref{ch:solvers}
returns the same four edges with no violations, and the exact MILP of
Chapter~\ref{ch:exact} reaches the same soft energy. Three independent
routes agree on one answer.

\section{Verification of the Encoding}
\label{sec:qubo:verify}

Two further checks were run on random instances.

\begin{enumerate}[label=\arabic*.]
\item \textbf{Exhaustive optimum check.} Over 38 randomly generated
instances, the exact bit-level optimum of Eq.~\eqref{eq:qubo:energy} never
contains a violated constraint, as the dominance argument predicts.
\item \textbf{Independent re-derivation.} The energy was re-derived from the
specification by separate code that evaluates every term directly, without
building $Q$. On 40 instances with 8 random states each, the matrix energy
agrees with the independent evaluation to within $2.3\times10^{-10}$.
\end{enumerate}

This is strong empirical validation of the implementation, not a formal
proof, and it is reported as such. It does, however, settle the attribution
question that dominates the next chapter: when a solver later returns a poor
answer, the fault lies with the solver and not with the encoding.

\section{Problem Size}

The flow family dominates the variable count, and the matrix is stored
densely, so memory grows with the square of the number of bits. On
range-limited random fleets with connectivity enforced, ten rovers give 427
binary variables (a 1.5~MB matrix) and twenty rovers 1076 variables (9.3~MB);
without connectivity, fifty rovers need only 328. The dense formulation
therefore targets tens of rovers rather than hundreds, a limit revisited in
Chapter~\ref{ch:conclusion}.
